\section{Worked examples: why the stronger interface matters}
\label{sec:worked}

The first three examples unpack the common-base case: selecting random
bits, preserving the output law, and composing an error budget. The
remaining examples return to boundary locality, trace limits, and named
coordinates. Each isolates the information needed by a later operation
without the surrounding formal development.

\subsection{Injective tape positions preserve a uniform experiment}
\label{sec:tape-example}

\begin{samepage}
\begin{proposition}\label{prop:tape}
Let $p:\{0,\ldots,m-1\}\to\{0,\ldots,B-1\}$ be injective. If
$\beta$ is uniform on $\{0,1\}^B$, then
$(\beta(p(0)),\ldots,\beta(p(m-1)))$ is uniform on $\{0,1\}^m$.
Every word has exactly $2^{B-m}$ preimages.
\end{proposition}
\begin{proof}
Prescribing the $m$ selected bits leaves precisely $B-m$ coordinates
free. Each assignment to these remaining coordinates determines one
preimage. Thus every word has the same number of preimages, and its
probability is $2^{B-m}/2^B=2^{-m}$.
\end{proof}
\end{samepage}

For the common-base machine, the relevant map is
$p(i)=s+3(m-1-i)+1$. Reversal and multiplication by three preserve
distinctness. For $i<m$, its value is less than $B$ whenever
$s+3m+1\leq B$. The latter bound is stronger than needed for the selected
positions because it also accounts for the surrounding execution.
The review's Lean file proves these arithmetic facts for arbitrary natural
$s,m,B$, including the empty-word case.

\Needspace{7\baselineskip}
The index map is defined by
\begin{lstlisting}
def physicalIndex (start m : Nat) (i : Fin m) : Nat :=
  start + 3 * (m - 1 - i.val) + 1
\end{lstlisting}
The proof uses the bounds stored in \src{Fin m} to justify subtraction
arithmetic. This matters: natural-number subtraction is truncated, so
injectivity would fail on unrestricted natural indices. A range type is
part of the mathematical interface, not a cosmetic annotation.

For comparison, if $p(0)=p(1)$, the words $01$ and $10$ have no
preimages. Reading fair individual bits does not by itself produce an
independent word. Proposition~\ref{prop:tape} explains why the source proves
injectivity of the physical positions before transporting the success
probability. The Python companion exhaustively checks 15 small choices
of $(s,m,B)$ against the exact fiber count.

\subsection{Observables determine a finite law; its mean does not}
\label{sec:law-example}

\begin{samepage}
\begin{proposition}\label{prop:law}
Let $\mu,\nu$ be probability measures on a finite set $X$. Then
$\E_\mu\varphi=\E_\nu\varphi$ for every $\varphi:X\to\R$ if and only
if $\mu=\nu$.
\end{proposition}
\begin{proof}
For each $x\in X$, take $\varphi=\mathbf1_{\{x\}}$ to obtain
$\mu(\{x\})=\nu(\{x\})$. Summing singleton masses gives equality on every
subset. The converse follows from the definition of finite expectation.
\end{proof}
\end{samepage}

In a uniform finite experiment, these expectations are normalized sums
with integer multiplicities. The Lean companion proves the corresponding
count statement: equality of unnormalized sums against every natural-valued
observable determines each multiplicity. If two sampling spaces have
different cardinalities $N,N'>0$, equality of their output laws instead
means $N'c(x)=Nc'(x)$ for each output $x$. The companion also proves this
cross-multiplied statement from equality of normalized observables.
For example, multiplicities $(1,1,1)$ and $(2,2,2)$ describe the same law,
although their unnormalized counts differ. It then checks the measures on
$X=\{0,1,2\}$ with multiplicities
\[
  (1,1,1)\quad\text{and}\quad(0,3,0).
\]
Both have total mass three and mean one after normalization. Their
probabilities of $\{0\}$ are $1/3$ and zero. Thus preservation of the
expected numerical output would be insufficient for an accuracy event.
The universal observable certificate in Section~\ref{sec:common-bases}
supports exactly the stronger conclusion required: substitute the
indicator of the desired relative-error interval.

This interpretation also explains why the source allows reindexing of
draw outcomes. A permutation can change the result returned for a
particular numbered bank without changing its distribution under uniform
sampling. Law equality is the appropriate interface between those
experiments; the deterministic compiler layer separately preserves the
execution result for a supplied bank.

\subsection{From a bounded experiment to a finite success count}
\label{sec:success-budget}

The preceding two interfaces can be composed with a quantitative error
bound. This gives a concrete way to read the final accuracy statement of
Section~\ref{sec:common-bases}.

\begin{samepage}
\begin{proposition}\label{prop:success-budget}
Let $X$ be finite, adjoin a failure symbol $\bot$, and let $U$ be an
ideal $X$-valued output. Let $V$ be a bounded implementation driven by a
uniform $m$-bit word, taking values in $X\cup\{\bot\}$.
For $\varphi:X\to[0,1]$, set
$\varphi(\bot)=0$. Suppose
\[
  \E\varphi(V)\geq\E\varphi(U)-\eta
  \quad\text{for every such }\varphi.
\]
Let $\eta,\delta_0\geq0$ and
$A\subseteq X$ satisfy $\Pr(U\in A)\geq1-\delta_0$.
If a machine obtains its word from fixed injective positions in a
uniform $B$-bit tape and returns $V$ on that word, then
\[
  \#\{\text{successful }B\text{-bit tapes}\}
  \geq(1-\delta_0-\eta)2^B,
\]
where success means returning an element of $A$.
\end{proposition}
\begin{proof}
Choose $\varphi=\mathbf1_A$ to obtain
$\Pr(V\in A)\geq1-\delta_0-\eta$. Proposition~\ref{prop:tape}
shows that the tape projection has the required uniform word law.
Multiplying its success probability by $2^B$ gives the count.
\end{proof}
\end{samepage}

For a single accuracy event, only its indicator is needed; the
all-observables hypothesis makes the estimate reusable for other events.
In the source's bounded draw-tree implementation, the fit condition
requires at most $M$ successive draws and $H\geq\lceil\log_2 a\rceil$
trial bits for every draw arity $a$. With $D$ trials per draw,
$\eta=M2^{-D}$, with failed runs assigned
observable value zero.\footnote{See
\href{https://github.com/openai/math/blob/fd4aeeb2ee4fc729c18d98444fed42fd0529eeeb/lean/OAI/Combinatorics/MatroidCounting/CommonBases.lean\#L9346-L9349}{\src{CommonBases.lean}, \src{run_lower}}.}
For a desired failure budget $0<\delta<1$, choosing
$\delta_0\leq\delta/2$ and $M2^{-D}\leq\delta/2$ suffices.
When $M>0$, one may take
$D\geq\max\{0,\lceil\log_2(2M/\delta)\rceil\}$.
The required tape length and execution bound must still be proved
separately. Thus the accuracy budget, the uniformity of the gathered
bits, and the deterministic machine specification are distinct inputs
to a single final guarantee.

A tight finite example joins the layers. Let the ideal probabilities of
$a,b$ be $6/8,2/8$, and assign the eight implementation words to five
copies of $a$, two of $b$, and one failure. The observable loss is exactly
$\varphi(a)/8$, hence at most $\eta=1/8$. For $A=\{a\}$,
$\delta_0=1/4$ and the lower bound is $5/8$. Embedding the three-bit word
at positions $7,4,1$ in a ten-bit tape gives exactly
$5\cdot2^7=640$ successful tapes out of $1024$. The Python companion
enumerates those tapes and all event indicators. This fixture illustrates
the composed bound independently of the upstream rejection sampler.

\subsection{Boundary values and ambient derivatives}
\label{sec:local-forms}

The Mahler case uses local equality to identify ambient derivatives.
On $\R^2$, the
one-forms $\omega=0$ and $\eta=(x^2+y^2-1)\,dx$ agree at every point
of the unit circle. Their ambient exterior derivatives need not agree:
$d\eta=-2y\,dx\wedge dy$, which is nonzero at $(0,1)$.
Equality on a neighborhood does imply equality of the ambient derivatives.
The Mahler cutoff argument uses this sufficient condition to transfer
its ambient derivative; agreement on tangent arguments suffices for the
later comparison of integration densities. Intrinsic differentiation
after pullback to the circle uses different data: both pulled-back forms,
and their exterior derivatives, vanish. Mathlib's
differential-form interface explicitly uses the derivative of a
form-valued field \citep{MathlibForms}.


\subsection{A cut-sensitive limit that is not pointwise}
\label{sec:trace-example}

Consider the commutative $C^*$-algebra $c_0$ of complex sequences tending
to zero, with pointwise operations and the supremum norm. On its positive
cone define
\[
  \tau(a)=\sum_{j=1}^{\infty}a_j\in[0,\infty].
\]
The trace is additive, positively homogeneous, and tracial. It is lower
semicontinuous because it is the supremum of its continuous finite partial
sums. Finite-support sequences belong to $\mathcal N_\tau$, so the closed
finite ideal $I_\tau$ in~\eqref{eq:finite-ideal} is all of $c_0$.

For every positive $a\in c_0$ and $\varepsilon>0$, the sequence
$(a-\varepsilon)_+$ has finite support. Consequently
\[
  \lim_{r\downarrow0}r\tau((a-\varepsilon)_+)=0,
  \qquad 0\leq\liminf_{r\downarrow0}r\tau(a).
\]
These are precisely the two inequalities in
\eqref{eq:trace-canonical-convergence} with limit trace
$\iota_{c_0}=0$. But for $a_j=1/j$ one has $\tau(a)=\infty$, hence
$r\tau(a)=\infty$ for every $r>0$. The pointwise limit at this element
is infinity. The finite ideal and canonical limit are nevertheless
unambiguous.

This classical example makes the reason for the positive cuts concrete
\citep{ElliottRobertSantiago2011}. Closure admits elements with infinite
trace, while every fixed positive cut is controlled by a finite value.
Replacing the source's convergence predicate by pointwise convergence
would therefore alter the theorem.


\subsection{Equal flags and different named coordinates}
\label{sec:named-maps}

The character-variety case preserves more than the flag of subspaces.
Two frames of a one-dimensional vector space can induce the
same flag but different named coordinate maps. If $e(u)=uv$ and
$f(u)=2uv$ over a field of characteristic zero, both give the flag
$0\subset\operatorname{span}(v)$, while
$e^{-1}(v)=1$ and $f^{-1}(v)=1/2$. The all-seams predicate must therefore
retain the named graded maps when those maps occur in the seam equation.
The coordinate equation consumes information that equality of the
underlying flag has forgotten. This is the small-scale reason for
retaining named graded maps throughout the surface construction.
